% !TEX program = xelatex
\documentclass[11pt,letterpaper]{article}

% RDT: Rasool Document Template v1.0
% Best result: compile with XeLaTeX or LuaLaTeX so fontspec can load
% Calibri. If Calibri is unavailable, the template falls back to
% Carlito. A pdfLaTeX fallback is included so Overleaf projects still
% compile if the wrong compiler is selected.

\usepackage[letterpaper,top=1in,bottom=1.5in,left=1.5in,right=1.5in]{geometry}
\usepackage{iftex}
\ifPDFTeX
  \usepackage[T1]{fontenc}
  \usepackage[utf8]{inputenc}
  \usepackage{tgheros}
  \renewcommand{\familydefault}{\sfdefault}
\else
  \usepackage{fontspec}
\fi
\usepackage{setspace}
\usepackage{titlesec}
\usepackage{caption}
\usepackage{booktabs}
\usepackage{enumitem}
\usepackage{array}
\usepackage{xcolor}
\usepackage[hidelinks]{hyperref}

\ifPDFTeX
\else
  \IfFontExistsTF{Calibri}{\setmainfont{Calibri}}{\setmainfont{Carlito}}
\fi

\definecolor{rdtnavy}{HTML}{0B2E59}
\setstretch{1.26}
\setlength{\parindent}{0pt}
\setlength{\parskip}{8.5pt}

\titleformat{\section}{\bfseries\MakeUppercase}{\thesection}{0.5em}{}
\titleformat{\subsection}{\bfseries}{\thesubsection}{0.5em}{}
\titleformat{\subsubsection}{\bfseries\itshape}{\thesubsubsection}{0.5em}{}
\titlespacing*{\section}{0pt}{11pt}{8.5pt}
\titlespacing*{\subsection}{0pt}{8.5pt}{8.5pt}
\titlespacing*{\subsubsection}{0pt}{8.5pt}{8.5pt}

\captionsetup{font={small,stretch=1.0},labelfont={bf,it},labelsep=endash,margin=1in}
\setlist[itemize]{leftmargin=0.5in,labelsep=0.25in}
\setlist[enumerate]{leftmargin=0.5in,labelsep=0.25in}

\newenvironment{rdtabstract}{\begin{quote}\noindent\textbf{\textit{Abstract}}---}{\end{quote}}
\newenvironment{rdtblockquote}{\begin{quote}}{\end{quote}}
\newcommand{\rdtrunin}[1]{\textbf{\textit{#1---}}}

\begin{document}

\begin{center}
{\color{rdtnavy}\fontsize{17}{19.55}\selectfont Rasool Document Template v1.0\\
For Use in FIN 4950\par}
\vspace{5.5pt}
{\color{rdtnavy}\rule{\linewidth}{0.4pt}}
\vspace{5.5pt}
Author Name\\
Southeastern Louisiana University\\
email@example.edu
\end{center}

\begin{rdtabstract}
RDT standardizes formatting for FIN 4950 documents while keeping
submissions readable and consistent. Abstracts are optional; use one
for longer assignments that need a brief summary.
\end{rdtabstract}

\section{Typography}

RDT keeps FIN 4950 submissions readable and consistent. Use the styles
in this template instead of formatting each paragraph by hand.

\subsection{Body text}

All main text in RDT uses Calibri regular at 11 points when available.
Body paragraphs use 1.26 line spacing, 8.5 points of spacing after each
paragraph, justified alignment, and no first-line indentation.

Use the built-in styles for headings, captions, body text, and
references.

\textbf{Bold} and \textit{italics} may be used for emphasis. Hyperlinks,
inline code such as \texttt{risk\_premium}, superscripts such as
\(x^2\), and subscripts such as \(\beta_1\) are allowed when they help
explain the work.

\subsection{Title and subtitle}

The title uses Calibri regular at 17 points, centered at the top of the
first page with 1.15 line spacing. It may span up to three lines when a
project title is unusually long. For this course template, the second
line must read ``For Use in FIN 4950.''

The author name, affiliation, and email lines appear below the title
block in body-size text and are centered. The affiliation line should
read ``Southeastern Louisiana University.'' The email line may be
omitted for anonymous submission or different instructor directions.

\subsection{Abstract}

When included, the abstract appears below the title block. The word
``Abstract'' is bold italic and followed by an em dash. The paragraph
uses side indents, body-size type, justified alignment, and the same
paragraph rhythm as the main text.

Use an abstract for longer assignments that need a brief summary. Short
assignments usually do not need one unless the prompt requires it.

\subsection{Headings}

Use numbered headings. Heading 1 is all caps. Heading 2 and Heading 3
use sentence case; Heading 2 is bold upright, and Heading 3 is bold
italic. Heading 4 is a run-in sidehead: bold italic text followed by an
em dash.

Headings should be clear and usually fit on one line. Word and LaTeX can
handle numbering through styles or commands; Google Docs and Pages may
require manual numbering or careful style setup.

\subsubsection{Heading 1}

Heading 1 is for major sections. It is all caps, with 11 points of
spacing before and 8.5 points after. Use it sparingly for the main
structure of the document. Examples include INTRODUCTION, DATA,
METHODS, RESULTS, DISCUSSION, REFERENCES, or APPENDICES.

\subsubsection{Headings 2 through 4}

Heading 2 is for major subsections inside a Heading 1 section. Heading
3 is for deeper explanation under Heading 2. Keep both in sentence case
instead of title case. Use Heading 4 as a labeled paragraph inside dense
material.

\rdtrunin{Heading 4 example} This sentence demonstrates a compact run-in
sidehead for a local label inside a paragraph.

\subsection{Page layout}

RDT uses US Letter paper, 8.5 by 11 inches. The top margin is 1 inch,
and the left, right, and bottom margins are 1.5 inches.

The page number should appear in the footer, centered or in a standard
footer position, about 1 inch from the bottom. Do not place decorative
elements in the margins. Keep body text, figures, tables, captions,
references, and appendices within the margins. Do not shrink margins,
font size, or spacing to fit extra material into a page-limited
assignment.

\section{Presentational elements}

Use figures, tables, equations, screenshots, charts, and other visuals
to clarify the text, not to avoid explaining it.

\subsection{Figures}

Figures should be centered and may use the full text block width when
needed, but they must remain inside the margins. Every figure should be
referenced in the body text and should have a descriptive caption.

If a figure has a white background, a thin border can improve
visibility. Figure captions go below figures and use small text with
clear side margins. The label, such as ``Figure 1,'' is bold italic and
followed by an em dash. Keep captions on the same page as the figure
whenever possible.

\begin{figure}[h]
\centering
\fbox{\parbox{0.72\linewidth}{\centering Insert figure, chart, map,
screenshot, equation, or model output here.}}
\caption{Example RDT figure caption placed below a figure and written as
a complete description.}
\end{figure}

\subsection{Tables}

Tables should be centered and fit within the margins. Table captions go
above tables. Figure captions go below figures; table captions go above
tables. Table captions follow the same format as figure captions.

Text in tables is generally left-aligned. Numeric columns are generally
right-aligned or decimal-aligned. Dense tables may use smaller text,
often matching caption size. Move large tables to appendices if they
interrupt the document. Use editable tables instead of screenshots.

\begin{table}[h]
\caption{Example RDT table with text, numeric values, and concise notes.}
\centering
\begin{tabular}{l r l}
\toprule
Measure & Value & Use in analysis \\
\midrule
Expected return & 0.082 & Forecasted annual return \\
Volatility & 0.165 & Risk comparison \\
Sample size & 120 & Number of observations \\
R-squared & 0.410 & Model fit summary \\
\bottomrule
\end{tabular}
\end{table}

\subsection{Additional elements}

RDT also supports quotes, lists, equations, code snippets, footnotes,
and appendices.

\subsubsection{Quotes}

Short quotes can remain inline with quotation marks and citations.
Set quotes longer than three lines as block quotes with 0.5-inch side
indents. Use brackets only to clarify replaced words. Prefer paraphrase
when a quote adds length without adding value.

\begin{rdtblockquote}
A clear financial analysis does more than report a number; it explains
what the number measures, why the measure is appropriate, and how the
result should guide the reader's decision.
\end{rdtblockquote}

\subsubsection{Lists}

Use bulleted lists for parallel items. Use numbered lists for sequence,
priority, or steps. Lists should use a 0.5-inch left indent, with the
bullet or number hanging by about 0.25 inch. Avoid long nested lists
unless they are necessary, and keep list items parallel.

\begin{itemize}
\item Use a bullet when the order of items does not matter.
\item Use parallel wording so list items scan cleanly.
\end{itemize}

\begin{enumerate}
\item Define the question.
\item Present the evidence.
\item Explain the implication.
\end{enumerate}

\subsubsection{Code, formulas, and technical notation}

Inline code may be used for variable names, file names, commands, or
functions. Longer code blocks may be separated from the body text.
Mathematical notation should be readable and consistently formatted.
Define variables before using them heavily. Avoid unexplained equations
or code blocks.

\section{Procedural elements}

This section explains citations, references, and appendices.

\subsection{In-text citations}

RDT uses APA-style in-text citations by default unless the assignment,
journal, client, or instructor requires another style. Citations should
appear close to the claim they support. Do not place one citation only
at the end of a long paragraph if several claims require support.

Examples include a parenthetical citation such as (Rasool, 2026), a
narrative citation such as Rasool and Smith (2026), multiple sources
such as (Rasool, 2026a; Rasool, 2026b), and a multi-author citation
such as Rasool et al. (2026). Citations should support specific claims,
data, methods, figures, and non-original ideas.

Footnotes may be used for brief clarification, but they are not
the default citation method in RDT. Footnotes should use small text with
the same line-spacing rhythm.

\subsection{Reference lists}

References appear in a dedicated References section. Only works cited
in the body should appear in the reference list. References should be
alphabetized by the first author's last name. If numbered references
are used, numbering should follow the alphabetized order.

Multiple works by the same author should be ordered chronologically.
Same-author, same-year works should use letters such as 2026a and
2026b. Reference lists generally do not count against page limits unless
the assignment says otherwise.

\section{References}

\begin{enumerate}
\item Rasool, A. (2026). \textit{Example title for a research report}.
Example Journal, 1(1), 1--10.
\item Smith, J., \& Rasool, A. (2026). \textit{Example title for an
applied policy brief}. Example Policy Review, 2(1), 15--25.
\end{enumerate}

\section{Appendices}

Appendices are optional and appear after References. If there are multiple appendices,
use a Heading 1 called ``APPENDICES'' and Heading 2 labels for each
appendix. If there is only one appendix, use ``Appendix: Descriptive
title.''

Use appendices for raw data, long tables, model diagnostics, extended
calculations, or extra figures. Do not put the main argument in an
appendix. The body should summarize appendix content and refer readers
to the relevant appendix. Appendices generally do not count against page
limits unless the instructor says otherwise.

\subsection{Appendix A: Example supplemental table}

This placeholder shows where optional supporting material may appear.
The main body should remain understandable without the appendix.

\end{document}
